% =====================================================================
%  GMN3 §3.4--3.6  带框穿墙 / motivic 穿墙公式 / 线算子的形变环
%  原文：Lamination-Paper.tex 1507--2020 行；PDF p.22--27（印刷 p.21--26）
%  编号公式 (3.25)--(3.63) 共 39 条（含脚注中 2 条），全部在本块内，故 \eqref 可用；
%  插图 1 幅（fig:zeta-u-wall）；脚注 3 条；跨块引用暂写原文编号并加 % TODO
% =====================================================================

\subsection{带框穿墙}\label{subsec:Halo-Wall-Crossing}

在第 3.2 节% TODO 合稿改 \ref{subsec:BPS-Deg-PSC}
中，我们定义了保护自旋特征标（PSC）$\fro(L,\gamma;u,\zeta)$，它「计数」带框 BPS 态。但正如我们在第 3.1 节% TODO 合稿改 \ref{sec:framed-BPS-states}
中所指出，当 $(\zeta,u)$ 越过一道 BPS 墙时，带框 BPS 态会与连续谱混合。因此，关于 $\fro(L,\gamma;u,\zeta)$ 不变性的标准论证失效了。这就是「穿墙」问题的带框版本。所幸它比「原味」BPS 态的完整穿墙问题要容易：从带框 BPS 谱中消失的态，正是我们在第 3.3 节% TODO 合稿改 \ref{subsec:Halos}
中描述过的那些简单的 halo 态。于是我们可以完全描述 $\fro(L,\gamma;u,\zeta)$ 在越过墙时的跳变。

考虑一条路径 $(u_t,\zeta_t)\in\widehat{\CB}\times\widehat{U(1)}$，它穿过一道墙 $\widehat{W}(\gamma_h)$。如上所述，当抵达该墙时 $r_{\rm halo}$ 趋于无穷。因此，一整个由 halo 束缚态构成的 Fock 空间——由电荷正比于 $\gamma_h$（比例系数为正）的 halo 粒子搭建而成——要么出现在 BPS 谱中，要么从其中消失。

现在让我们描述这件事对（带框）保护自旋特征标的影响。设 $\{x_\gamma\}$ 是 $\Gamma$ 的群代数的一组基，使得
\begin{equation} \label{eq:untwisted}
x_{\gamma} x_{\gamma'} = x_{\gamma +
\gamma'}.
\end{equation}
然后构造生成泛函
\begin{equation}\label{eq:gen}
F(u, L, \zeta, \{x_\gamma\}; y) := \sum_{\gamma} \fro(u, L, \zeta,
\gamma; y) x_\gamma.
\end{equation}
（为了避免记号变得过于笨重，当某些变量可以安全地当作隐含时，我们常常会把它们省略。）

上文所述的每个类型为 $m, m', \gamma_h$ 的产生算子，都为 Fock 空间上迹的乘积表示贡献一个因子
\begin{equation}
(1 + (-1)^{m} y^{m+m'} x_{\gamma_h} )^{  a_{m,\gamma_h}}
\end{equation}
定义 $\Gamma_h := \IZ \gamma_h$，并对 $\gamma_c \in \Gamma_{\CL}$ 令
\begin{equation}
F_{\bar \gamma_c} := \sum_{\gamma'\in \Gamma_h} \fro(u, L, \zeta,
\gamma_c +\gamma'; y) x_{\gamma_c+\gamma'}.
\end{equation}
（这个求和只依赖于投影 $\bar\gamma_c\in \Gamma_L/\Gamma_h$。）\zy{原文把 $\gamma_c$ 所取的格写作 $\Gamma_{\mathcal L}$，紧接的投影又写 $\Gamma_L$，二者应指同一个格（线算子的电荷格）。}
这是 $F$ 中对应于电荷类型为 $\gamma_c + \ell \gamma_h$ 的态的那一部分。我们用 $F_{\bar\gamma_c}^\pm$ 分别记这个生成函数在墙的两侧——即分别满足 ${\rm Im}(Z_{\gamma_h}(u)/\zeta)>0$ 与 ${\rm Im}(Z_{\gamma_h}(u)/\zeta)<0$ 的一侧——的取值。于是有
\begin{equation} \label{eq:transform}
F^\pm_{\bar\gamma_c} = F^\mp_{\bar\gamma_c}  \prod_{\gamma_h}
\prod_{m=-M_{\gamma_h}}^{M_{\gamma_h}} \prod_{m'= - 2J_{\gamma_c,
\gamma_h}}^{2J_{\gamma_c, \gamma_h}} (1 + (-1)^{m } y^{m+m'}
x_{\gamma_h} )^{  a_{m,\gamma_h}}.
\end{equation}
这里我们稍微滥用了记号：对 $\gamma_h$ 的第一个乘积，指的是对所有给出同一道墙 $\widehat{W}(\gamma_h)$ 的 halo 电荷取乘积。左端究竟取 $F^+$ 还是 $F^-$（也就是我们究竟是获得还是失去一个 Fock 空间），取决于穿墙的方向以及 $\langle \gamma_c,\gamma_h\rangle$ 的符号。（更精确的陈述见下文式 \eqref{eq:spin-ks-tmn-ii}。）

这个变换看起来像是形式变量 $x_\gamma$ 的一次「坐标变换」的效果：它把每个 $x_{\gamma_c}$ 乘上 $x_{\gamma_h}$ 的一个有理函数。遗憾的是，这种解释与乘法规则 \eqref{eq:untwisted} 并不完全相容。但当 $y=1$ 时它确实行得通，下面我们就来说明。此时 \eqref{eq:transform} 特化为
\begin{equation} \label{eq:tspec}
F^\pm_{\bar\gamma_c}  = F^\mp_{\bar\gamma_c}  \prod_{\gamma_h}
\prod_{m=-M_{\gamma_h}}^{M_{\gamma_h}} (1 + (-1)^{m}
x_{\gamma_h} )^{  \lvert \langle \gamma_h, \gamma_c\rangle \rvert a_{m,\gamma_h} }.
\end{equation}
而「坐标」变换
\begin{equation}\label{eq:y-one-prod}
x_\gamma \to x_\gamma  \prod_{\gamma_h}
\prod_{m=-M_{\gamma_h}}^{M_{\gamma_h}}  (1 + (-1)^{m } x_{\gamma_h} )^{\langle \gamma_h, \gamma\rangle a_{m,\gamma_h}}
\end{equation}
在我们从 ${\rm Im}(Z_{\gamma_h}(u)/\zeta) < 0$ 的一侧穿墙到 ${\rm Im}(Z_{\gamma_h}(u)/\zeta) > 0$ 的一侧时，恰好重现 \eqref{eq:tspec}。此外，它与乘法法则 \eqref{eq:untwisted} 相容。

若 $y=-1$，\eqref{eq:transform} 中的因子 $y^{m'}$ 与 $- (-1)^{\langle \gamma_h, \gamma_c\rangle }$ 相同。有一个简单的技巧可以驯服这些多出来的符号：把乘法规则从 \eqref{eq:untwisted} 改成
\begin{equation} \label{eq:twisted-mult}
\hat x_{\gamma} \hat x_{\gamma'} = (-1)^{\langle \gamma, \gamma'\rangle } \hat x_{\gamma +\gamma'}.
\end{equation}
于是，利用关系
$\hat x_{\gamma_c + n \gamma_h} = \hat x_{\gamma_c} \left( (-1)^{\inprod{\gamma_h, \gamma_c}} \hat x_{\gamma_h} \right)^n$，
我们得到 \eqref{eq:tspec} 的一个类比
\begin{equation} \label{eq:transform2}
F^\pm_{\gamma_c}  = F^\mp_{\gamma_c}  \prod_{\gamma_h}
\prod_{m=-M_{\gamma_h}}^{M_{\gamma_h}} (1 -
\hat x_{\gamma_h} )^{ \lvert \inprod{\gamma_h, \gamma_c} \rvert a_{m,\gamma_h}}.
\end{equation}
变换规则
\begin{equation}\label{eq:y-min-one-gen}
\hat x_\gamma \to \hat x_\gamma  \prod_{\gamma_h}
(1 - \hat x_{\gamma_h} )^{\inprod{\gamma_h, \gamma} \Omega(\gamma_h)}
\end{equation}
作用于 $y=-1$ 的生成泛函时，正确地重现 \eqref{eq:transform2}，并且与扭乘法规则 \eqref{eq:twisted-mult} 相容。

更仔细地陈述如下：若 $\zeta_{\rm cw}$ 位于 BPS 墙 $\widehat{W}(\gamma_h)$ 的顺时针一侧，$\zeta_{\rm ccw}$ 位于逆时针一侧，则在略去所有其他无关指标后，生成函数 $F$ 满足
\begin{equation}\label{eq:Careful-rule}
F\left(\zeta_{\rm cw}, \hat x_\gamma\right) = F\left(\zeta_{\rm ccw} , \hat x_\gamma (1- \hat x_{\gamma_h})^{-\langle \gamma, \gamma_h\rangle \Omega(\gamma_h)} \right).
\end{equation}

上面用到的变换 \eqref{eq:y-min-one-gen}，正是 \cite{ks1} 中引入的辛同胚：定义 $\CK_{\gamma_h}$ 为
\begin{equation}\label{eq:KS-sympo}
\CK_{\gamma_h}(\hat x_\gamma) := \hat x_\gamma (1-\hat x_{\gamma_h})^{\langle
\gamma, \gamma_h \rangle },
\end{equation}
则 \eqref{eq:y-min-one-gen} 就简单地是
$\hat x_\gamma \to \CK_{\gamma_h}^{-\Omega(\gamma_h)}(\hat x_\gamma)$。

\subsubsection{非对易变量与量子双对数}

对一般的 $y$，我们无法把 \eqref{eq:transform} 解释为一次变量替换——至少无法用\emph{对易}变量做到。但假设我们引入满足下列关系的形式变量 $X_\gamma$
\begin{equation}\label{eq:Heis-Alg}
X_\gamma X_{\gamma'} = y^{\inprod{\gamma, \gamma'}} X_{\gamma +
\gamma'},
\end{equation}
并再次考虑生成函数
\begin{equation}\label{eq:form-gen}
F(u,L,\zeta,\{X_\gamma\};y) := \sum_{\gamma} \fro(u, L, \zeta,
\gamma;y) X_\gamma,
\end{equation}
这一次把它看成这些\emph{非对易}变量的函数。在这种情况下，正如我们马上会看到的，增加或减少 halo 束缚态的效果，可以漂亮地概括为 $X_\gamma$ 的某个变换。这个变换由与 Faddeev 和 Kashaev \cite{Faddeev:1993rs} 的量子双对数作共轭来实现。（这个函数所享有的各种性质与别名，在 \cite{Teschner:2005bz} 的附录 A 中有很好的总结。）

首先我们需要一条数学引理。设 $\langle \gamma_c, \gamma_h \rangle = n$ 不为零。我们在群代数中展开对易变量，以此定义 $y$ 的洛朗多项式：
\begin{equation}
x_{\gamma_c} (1 + y^{n-1} x_{\gamma_h})(1 + y^{n-3}
x_{\gamma_h})\cdots (1 + y^{3-n} x_{\gamma_h})(1 + y^{1-n}
x_{\gamma_h}) = \sum_{j=0}^{\vert n\vert}  P^{(n)}_j(y) x_{\gamma_c
+ j \gamma_h}
\end{equation}
从物理上看很显然，$P^{(n)}_j(y)$ 就是第 $j$ 个反对称积 $\Lambda^j \rho_{\vert n \vert }$ 的特征标，其中 $\rho_{N}$ 是 $SU(2)$ 的 $N$ 维不可约表示。\footnote{顺便说一句，有一个「$q$-二项式定理」也给出同样的识别（取 $n>0$），这既有趣又会很有用：
\begin{equation}
P^{(n)}_j(y) = \frac{ [n]_y ! }{ [j]_y ! [n-j]_y !}:=\genfrac{[}{]}{0pt}{}{n}{j}
\end{equation}
其中
\begin{equation}
[n]_y := \frac{y^{n} - y^{-n}}{y - y^{-1}} .
\end{equation}
}

我们的引理断言：在展开非对易表达式
\begin{equation}
X_{\gamma_c} \Phi_n( X_{\gamma_h}) = \sum_{j=0}^{\vert n \vert}
P^{(n)}_j(y) X_{\gamma_c + j \gamma_h}
\end{equation}
时出现的，是\emph{同一组}洛朗多项式，其中
\begin{equation}
\Phi_n(\xi) := \begin{cases}  \prod_{s=1}^{n} (1 + y^{-(2s-1)} \xi )
& n >0 \\
 1 & n=0 \\
 \prod_{s=1}^{\vert n\vert } (1 + y^{(2s-1)} \xi ) & n< 0
 \end{cases}
\end{equation}
这条引理的证明是直接的，从略。

利用这条引理，我们可以断言：穿墙的效果，是把 \eqref{eq:form-gen} 中\emph{所有}非对易变量作如下变换
\begin{equation}\label{eq:Spin-KS-tmn}
X_\gamma \to X_\gamma \prod_{m=-M_{\gamma_h}}^{M_{\gamma_h}}
\Phi_{\langle \gamma, \gamma_h\rangle} ( (-1)^m y^m
X_{\gamma_h})^{\epsilon a_{m,\gamma_h}}
\end{equation}
在这个变换律中，我们正在穿过墙 $\widehat{W}(\gamma_h)$。若我们从 $\langle \gamma_h, \gamma \rangle \ {\rm Im}(Z_{\gamma_h}(u)/\zeta) < 0$ 的一侧穿到 $\langle \gamma_h, \gamma \rangle \ {\rm Im}(Z_{\gamma_h}(u)/\zeta) > 0$ 的一侧，则符号 $\epsilon=+1$，因为此时我们在核心电荷 $\gamma$ 周围\emph{获得}了一个 halo 粒子的 Fock 空间。若沿相反方向穿墙，则 $\epsilon=-1$，因为此时我们\emph{失去}了一个 halo 粒子束缚态的 Fock 空间。若用 $F^+, F^-$ 分别记形式生成函数 \eqref{eq:form-gen} 在 ${\rm Im}(Z_{\gamma_h}(u)/\zeta) > 0$ 与 $<0$ 一侧的取值，则（略去所有其他无关变量）：
\begin{equation}\label{eq:spin-ks-tmn-ii}
F^+(X_\gamma) = F^-\left( X_\gamma \left( \prod_{m=-M_{\gamma_h}}^{M_{\gamma_h}}
\Phi_{\langle \gamma, \gamma_h\rangle} ( (-1)^m y^m
X_{\gamma_h})^{a_{m,\gamma_h}}\right)^{{\rm sign}(\langle \gamma_h,
\gamma\rangle ) }\right)
\end{equation}

变换 \eqref{eq:Spin-KS-tmn} 可以用量子双对数的共轭来优雅地概括。定义
\begin{equation}
\Phi(X) := \prod_{k=1}^\infty ( 1 + y^{2k-1} X)^{-1}
\end{equation}
那么变换 \eqref{eq:Spin-KS-tmn} 与 \eqref{eq:spin-ks-tmn-ii} 就简单地等价于下述规则：
\begin{equation}\label{eq:Spin-KS-tmn-S}
F^+(X_\gamma) =   S_{\gamma_h} F^-(X_\gamma ) S_{\gamma_h}^{-1}
\end{equation}
其中
\begin{equation}\label{eq:Spin-KS-ii}
S_{\gamma_h} =  \prod_{\gamma_h} \prod_{m=-M_{\gamma_h}}^{M_{\gamma_h}} \Phi (
(-1)^m y^m X_{\gamma_h})^{a_{m,\gamma_h}}
\end{equation}
这里第一个乘积同样表示对与 $\gamma_h$ 平行的所有电荷取乘积。相应地，带框简并度的生成函数 \eqref{eq:form-gen} 按与 $S_{\gamma_h}$ 的共轭而变换。

\subsection{motivic 穿墙公式}\label{subsec:Motivic-WCF}

设 $\widehat{\CB}$ 是库仑分支 $\CB$ 光滑部分的万有覆盖，并考虑
\begin{equation}\label{eq:Chamber-Space}
\Xi: =  \widehat{\CB} \times \widehat{\IC}^* - \cup_{\gamma: \exists u: \Omega(\gamma;u) \ne 0 }  \widehat{W}(\gamma)
\end{equation}
这个空间被分成若干腔——即 $\Xi$ 的连通分支。我们用指标 $c$ 来标记这些腔。考虑 $\Xi$ 的两个腔 $c_1,c_2$。考虑 $\widehat{\CB} \times \widehat{\IC}^*$ 中连接这两个腔的一条路径 $\CP$。位于这两个腔中的线算子 $L$ 的生成函数 \eqref{eq:form-gen}，将由形如 \eqref{eq:Spin-KS-tmn-S} 的变换的某个复合相联系：\footnote{请注意，我们在这里悄悄引入了一个假设：生成函数可以从单位圆 $\vert \zeta \vert =1$ 延拓到 $\IC^*$。这个假设在当前的论证中并不本质，在第 4 节与第 5 节的讨论中也不本质。但在第 6 节中它将很重要。}% TODO 合稿改 \ref{sec:Formal-line-operators}、\ref{sec:Cluster-Algebras}、\ref{sec:linehol}
\begin{equation}\label{eq:Path-Transf}
F(L;c_1) = S(\CP)  F(L;c_2)
\end{equation}
其中 $S(\CP)$ 是与路径 $\CP$ 从 $c_1$ 走到 $c_2$ 途中所穿过的各道墙 $\widehat{W}(\gamma_h^i)$ 相对应的变换的路径有序乘积
\begin{equation}
S(\CP)= \prod_{\gamma_h^i} {\rm Ad}(S_{\gamma_h^i} ).
\end{equation}
我们希望断言：变换之积 $\prod_{\gamma_h^i} {\rm Ad}(S_{\gamma_h^i} )$ 与连接 $c_1$ 到 $c_2$ 的路径 $\CP$ 无关。只要理论拥有足够多的线算子 $L$，使得知道变换 \eqref{eq:Path-Transf} 的作用就强到足以完全确定 $S(\CP)$，这一点便成立。我们并没有完全一般的理由说明它为何成立（事实上，我们也没有完全一般的理由说明为什么竟然会存在任何线算子）。不过，我们将在下面的例子中看到，线算子总是足够多；事实上，在这些例子里，库仑分支上任一点处的线算子与电磁荷之间都存在一一对应，并且可以把任意 $X_\gamma$ 还原为各函数 $F(L,\{X_\gamma\})$ 的线性组合。进一步的讨论见下文第 12 节% TODO 合稿改 \ref{sec:Tropic}
。

\insfig[0.8\textwidth]{zeta-u-wall}{从 $(\zeta_1, u_1)$ 到 $(\zeta_2, u_2)$ 的两条路径。BPS 射线在联合的 $(\zeta, u)$ 空间中以余维 1 的轨迹出现。每道 BPS 墙 $\widehat{W}(\gamma)$ 都附着一个坐标变换。如图所示，墙可以汇合与分裂，因此两条路径一般会穿过不同的 BPS 墙集合。尽管如此，相应坐标变换的复合仍必须与路径无关。}

$S(\CP)$ 的路径无关性，实际上是 Kontsevich 与 Soibelman \cite{ks1} 的「motivic 穿墙公式」的一个版本。为看出这一点，考虑 $\widehat{\CB} \times \widehat{\IC}^*$ 中连接 $(\zeta_1, u_1)$ 到 $(\zeta_2, u_2)$ 的两条路径，如图 \ref{fig:zeta-u-wall} 所示；并比较 $\Xi$ 中包含 $(\zeta_1, u_1)$ 与包含 $(\zeta_2,u_2)$ 的两个腔 $c_1,c_2$ 的生成函数 $F(L, c_1)$ 与 $F(L,c_2)$。假设我们沿路径 $(\zeta_1, u_1) \to (\zeta_2, u_1) \to (\zeta_2, u_2)$ 走；并假设当 $u$ 沿从 $u_1$ 到 $u_2$ 的路径变化时，没有任何 $\arg Z_\gamma(u)$ 进入或离开区间 $(\arg \zeta_1, \arg \zeta_2)$。那么所有函数 $F(L,c_1)$ 都按下述有序乘积变换：
\begin{equation}
{\bS}(\zeta_1, \zeta_2;u_1) = \prod_{\gamma: \arg Z(\gamma) \in (\zeta_1, \zeta_2)} {\rm Ad}(S_\gamma).
\end{equation}
另一方面，若我们改走路径 $(\zeta_1, u_1) \to (\zeta_1, u_2) \to (\zeta_2, u_2)$，则它们改为按 ${\bS}(\zeta_1, \zeta_2;u_2)$ 变换。于是，对所有线算子 $L$ 都有
\begin{equation}\label{eq:SF-SF}
{\bS}(\zeta_1, \zeta_2;u_1)F(L,c_1)={\bS}(\zeta_1, \zeta_2;u_2)F(L,c_1).
\end{equation}

我们希望由 \eqref{eq:SF-SF} 断定这两个变换确实相等：
\begin{equation} \label{eq:s-equal}
{\bS}(\zeta_1, \zeta_2;u_1) = {\bS}(\zeta_1, \zeta_2;u_2).
\end{equation}
正如上面讨论过的，我们必须假设存在足够多的线算子 $L$，才能推出 \eqref{eq:s-equal}。

与标准讨论的一点小差别是：我们在这里只得到了量子双对数的乘积在「伴随表示」中的等式。不过，按照从方程 \eqref{eq:y-one-prod} 到 \eqref{eq:y-min-one-gen} 的讨论可知，特化到 $y=-1$ 便给出等式 \eqref{eq:s-equal}，此时
\begin{equation}\label{eq:Symplectic-S}
{\bS}(\zeta_1, \zeta_2;u_1) = \prod_{\gamma: \arg Z(\gamma) \in (\zeta_1, \zeta_2)} \CK_\gamma^{\Omega(\gamma_\BPS; u_1)}.
\end{equation}
以 \eqref{eq:Symplectic-S} 给出的 ${\bS}$ 代入的方程 \eqref{eq:s-equal}，就是 \cite{Gaiotto:2008cd,Gaiotto:2009hg} 中所讨论的 Kontsevich--Soibelman 穿墙公式（WCF）的形式。知道所有 $\zeta_{1,2}$ 对应的 ${\bS}(\zeta_1, \zeta_2;u)$，等价于知道 $u$ 处的 BPS 指标集合。而要确定完整的 BPS 谱，还需要更多信息，即保护自旋特征标。

本节最后，我们对一些相关文献作些评论。涉及非对易（$q$-形变）量子双对数的穿墙公式出现在 Kontsevich 与 Soibelman 的工作 \cite{ks1} 中。经典极限（$q = -1$）也在那里讨论过，并直接给出了 $\N=2$ 理论中第二螺旋度超迹的穿墙公式，后者在 \cite{Gaiotto:2008cd} 中从物理上得到了证明。然而，起初并不完全清楚如何把这一成功推广到 $q$-形变的情形。文献 \cite{Dimofte:2009bv} 提出，应当把 $q$ 等同于计数 BPS 态自旋的参数 $y$；文献 \cite{Dimofte:2009tm} 把这个提议精确化为一个穿墙公式，并证明它在若干例子中成立。尽管如此，一般的物理证明仍然缺失。文献 \cite{Cecotti:2009uf} 给出了一个新的论证，把对易关系 \eqref{eq:Heis-Alg} 与穿墙公式联系到 A 模型开拓扑弦，进而联系到 Chern--Simons 理论。而我们在这里给出了一个直接证明，用的只是四维规范理论的一般概念。

\subsection{线算子的形变环}\label{subsec:Deformed-Ring}

正如我们在第 2.2 节% TODO 合稿改 \ref{subsubsec:RingStructure}
中提到过的，位于不同点 $\vec x_i\in \IR^3$、保持 $osp(4^*\vert 2)_\zeta$ 的平行超对称线算子可以相乘，并重新展开为线算子之和。这个运算定义了一个环：它应当是交换的，因为 $\IR^3$ 中没有自然的点序；它也应当是结合的，因为它们都是算子。自然要问：这个环结构与带框保护自旋特征标、带框 BPS 简并度之间有什么关系？结果表明，方程 \eqref{eq:form-gen} 中定义的生成函数 $F(L)$ 很适合讨论这个问题。在本节中我们将看到，可以用它们来定义线算子环的一个非常有意思的非交换形变。

假设我们把一个线算子 $L$ 放在 $\vec x=0$，另一个 $L'$ 放在沿 $z$ 轴平移的 $\vec x' = (0,0,z)$。注意由 $J_3$ 生成的旋转是这个位形的对称性，但 $so(3)$ 对称性已经破缺为 $so(2)$。若 $z\ne 0$，且 $L$ 与 $L'$ 按第 2.3 节% TODO 合稿改 \ref{subsec:LabelsForLines}
所解释的意义彼此互局域，则希尔伯特空间 $\CH_{L_{\vec x} L'_{\vec x'}}$ 是 $U(1)_{\vec x \vec x'}\subset SU(2)$ 的一个表示。若 $L$ 与 $L'$ 是单对象，则希尔伯特空间 $\CH_{LL'}$ 也按 $\Gamma + \gamma_L + \gamma_{L'}$ 分次，且分次中的各直和项满足
\begin{equation}\label{eq:LL-grade}
\CH_{L_{\vec x} L'_{\vec x'},\gamma_0}^{\rm BPS} = \bigoplus_{\gamma+ \gamma'=\gamma_0} \CH_{L,\gamma}^{\rm BPS}
 \otimes \CH_{L',\gamma'}^{\rm BPS}
 \otimes N_{\gamma,\gamma'}.
\end{equation}
这里 $N_{\gamma,\gamma'}$ 是 $U(1)_{\vec x \vec x'}$ 的一维表示。它来自这对双荷子所激发的电磁场。因此它对保护自旋特征标的贡献是\footnote{这一点应当与「原味」BPS 态的本原穿墙公式中的情形相对照。在那种情形，电磁场具有 $SU(2)$ 自旋 $2j = \vert \langle \gamma, \gamma' \rangle \vert -1$。多出来的这个 $-1$ 来自两个分离组分的质心自由度的对齐。}
\begin{equation}
{\Tr}_{N_{\gamma,\gamma'}} y^{2J_3} = y^{\langle \gamma, \gamma' \rangle}.
\end{equation}
回想生成函数 $F(L)$ 的定义（\eqref{eq:form-gen}）以及海森堡关系 \eqref{eq:Heis-Alg}，我们看到方程 \eqref{eq:LL-grade} 可以用来定义线算子的一个\emph{非交换}乘积：
\begin{equation}\label{eq:Non-Comm-Line-Prod}
F(L\circ_y L') := F(L) F(L').
\end{equation}
一般而言 $F(L)F(L') \ne F(L') F(L)$，从下面讨论的例子中可以看得很清楚。

在第 2.2 节% TODO 合稿改 \ref{subsubsec:RingStructure}
中，我们描述了线算子上的一个交换乘积。因此，发现这个乘积还有一个非交换形变，或许会令人意外。区别在于：是先取算子的 OPE $\lim_{z\to 0} LL'$，得到线算子之和；还是先在希尔伯特空间上取迹。这两种做法并不对易：
\begin{equation}
\lim_{z\to 0} F(LL')\ne F(\lim_{z\to 0} LL').
\end{equation}

为了更好地理解交换乘积与非交换乘积之间的关系，让我们假设存在一族单对象 $L_i$，它们生成全部线算子。例如，第 2.3 节% TODO 合稿改 \ref{subsec:LabelsForLines}
中讨论的任一族线算子 $\CL$ 都可以。单对象具有这样的性质：任何线算子都可以写成 $\sum c_i L_i$，其中 $c_i$ 是正整数。为定义环结构，只需计算融合系数
\begin{equation}\label{eq:simp-mult}
\lim_{z\to 0} L_i L_j = \sum_k c_{ij}^k L_k.
\end{equation}
然而，希尔伯特空间的类似分解取如下形式
\begin{equation}
\CH_{L_i L_j}= \oplus_k  N_{ij}^k \otimes \CH_{L_k}
\end{equation}
这是 $U(1)_{\vec x \vec x'}$ 表示的等式。令人意外的是（至少令作者们意外），在某些例子中 $N_{ij}^k$ 可以是 $U(1)_{\vec x \vec x'}$ 的非平凡表示。我们引入
\begin{equation}
c_{ij}^k(y) := {\Tr}_{N_{ij}^k} y^{2J_3}
\end{equation}
这是 $y$ 的一个洛朗多项式，系数为正整数，并满足 $c_{ij}^k(1) = c_{ij}^k$。用这些记号，线算子的交换环 \eqref{eq:simp-mult} 被形变为
\begin{equation}
F(L_i\circ_y L_j) = \sum_{k} c_{ij}^k(y) F(L_k)
\end{equation}
由于 $\CN=2$ 理论是宇称不变的，我们有 $c_{ij}^k(y)= c_{ji}^k(1/y)$。注意因此当 $y^2=1$ 时乘积 $\circ_y$ 是交换的。

为了理解为什么 $N_{ij}^k$ 必须是 $SO(2)$ 的非平凡表示，我们在认定两个分离线算子系统的正确量子数 $J_3$（从而 $F = 2 J_3 \,{\rm mod}\, 2$）时需要谨慎一些。首先，为了让 $J_3$ 能够被定义，我们需要把两个算子沿 $z$ 轴靠拢。此外，单个线算子的 $J_3$ 定义是明确的（常数平移会与非阿贝尔旋转群不相容）；但对含两个缺陷的系统，未破缺的旋转群是阿贝尔的，因此在比较 $LL'$ 的带框 PSC 定义与各个线算子 $L_i$ 的带框 PSC 定义中出现的 $J_3$ 时，可能会出现整体的常数平移。这种平移歧义由重数 $N_{ij}^k$ 的 $SO(2)$ 分次所体现。

注意，我们的乘积公式通过了一项非常严格的检验，即与穿墙的相容性：只要乘积按规则 $X_\gamma X_{\gamma'} = y^{\langle \gamma,\gamma' \rangle} X_{\gamma+\gamma'}$ 进行，穿墙变换的作用就与乘积可交换。
